%! AP Statistics — Unit 1, Lesson 1.10 — Homework KEY
\documentclass[10pt]{article}
\usepackage{apstats-article}
\usepackage{apstats-key}

%=============================================================================
\begin{document}

\pageheader{Unit 1, Lesson 1.10}{Homework — Key}
\begin{tabularx}{\linewidth}{X X X}
  Name:\;\ans{Key} & Date:\; & Period:\;
\end{tabularx}

\vspace{0.10in}

\begin{blurbbox}[Part A \quad The Empirical Rule — Key]{navylight}
\small IQ scores: $N(100, 15)$.

\begin{enumerate}[leftmargin=1.5em, itemsep=6pt]

  \item Axis values (left to right): \ans{55, 70, 85, 100, 115, 130, 145}.

  \item Between 85 and 115 ($\mu \pm 1\sigma$): \ans{68\%}\quad
        Between 70 and 130 ($\mu \pm 2\sigma$): \ans{95\%}

  \item Above 130 is above $\mu + 2\sigma$: \ansline{$(100\%-95\%)/2 = 2.5\%$. About 2.5\% of people score above 130.}

  \item $145 = \mu + 3\sigma$: \ansline{By the Empirical Rule, 99.7\% of scores fall between 55 and 145, so only about 0.15\% of people score above 145. A score of 145 is extremely unusual.}

\end{enumerate}
\end{blurbbox}

\vspace{0.10in}

\begin{blurbbox}[Part B \quad $z$-Scores and the Normal Table — Key]{greenacc}
\small Wait times: $N(4.5, 1.2)$ min.

\begin{enumerate}[leftmargin=1.5em, itemsep=6pt]

  \item $z = (6 - 4.5)/1.2 = 1.5/1.2 = \ans{+1.25}$\quad
        Interpretation: \ansline{A wait of 6 minutes is 1.25 standard deviations above the mean wait time.}

  \item $P(\text{wait}>6) = 1 - P(Z<\ans{1.25}) = 1 - \ans{0.8944} = \ans{0.1056}$\\
        About 10.6\% of customers wait more than 6 minutes.

  \item $z(3) = (3-4.5)/1.2 = -1.5/1.2 = \ans{-1.25}$\quad
        $z(6) = \ans{+1.25}$ (from above)\\[3pt]
        $P(3<\text{wait}<6) = P(Z<\ans{1.25}) - P(Z<\ans{-1.25}) = \ans{0.8944} - \ans{0.1056} = \ans{0.7888}$

  \item $x = 4.5 + (-1.28)(1.2) = 4.5 - 1.536 = \ans{2.96}$ minutes (approximately 3.0 min).

\end{enumerate}
\end{blurbbox}

\vspace{0.10in}

\begin{blurbbox}[Part C \quad AP-Style Free Response — Key]{redacc}
\small Graduate exam scores: $N(500, 100)$.

\begin{enumerate}[leftmargin=1.5em, itemsep=6pt]

  \item $z = (650-500)/100 = \ans{+1.50}$. $P(Z < 1.50) = 0.9332$.\\
        $P(\text{score}>650) = 1 - 0.9332 = \ans{0.0668}$.\\
        \ansline{About 6.7\% of test-takers scored above 650.}

  \item $z = (400-500)/100 = \ans{-1.00}$. $P(Z < -1.00) = \ans{0.1587}$.\\
        \ansline{About 15.9\% of test-takers scored below 400.}

  \item $z^* = 1.645$. $x = 500 + 1.645(100) = 500 + 164.5 = \ans{664.5} \approx 665$.\\
        \ansline{A score of approximately 665 separates the top 5\% from the rest.}

  \item \ansline{The requirement of 650 is quite selective. Only about 6.7\% of test-takers score at or above 650 (from part a). This means roughly 93\% of applicants would be excluded by this cutoff alone.}

\end{enumerate}
\end{blurbbox}

\begin{teachernote}
Part A Q4: $P(Z > 3) = 1 - 0.9987 = 0.0013$; one tail is $\approx 0.0013/2 \approx 0.00065 \approx 0.065\%$. The Empirical Rule gives 0.15\% for one tail beyond 3$\sigma$ — close enough.\\
Part B Q4: \texttt{invNorm}(0.10, 4.5, 1.2) $= 2.962$. Accept 2.96 or 3.0 min.\\
Part C Q1: Table A, row 1.5, col .00 $= 0.9332$. Q3: $z^* = 1.645$ is between $z=1.64$ (0.9495) and $z=1.65$ (0.9505).
\end{teachernote}

\end{document}
